\documentclass[conference]{IEEEtran}
\IEEEoverridecommandlockouts
\usepackage{cite}
\usepackage{amsmath,amssymb,amsfonts}
\usepackage{algorithmic}
\usepackage{graphicx}
\usepackage{textcomp}
\usepackage{xcolor}
\usepackage{multirow}
\usepackage[brazil]{babel}
\def\BibTeX{{\rm B\kern-.05em{\sc i\kern-.025em b}\kern-.08em
    T\kern-.1667em\lower.7ex\hbox{E}\kern-.125emX}}

\begin{document}

\title{An\'alise Comparativa de Arquiteturas de Deep Learning para
Classifica\c{c}\~ao de Retinopatia Diab\'etica}

\author{
\IEEEauthorblockN{
Gabriel de Sousa Nobre,
Victor Lins Gurgel do Amaral,
Lorenzo Barros Calheiros Pinheiro,
Davi Oliveira de Alencar}
\IEEEauthorblockA{
\textit{Gradua\c{c}\~ao em Ci\^encia da Computa\c{c}\~ao} \\
\textit{Universidade de Fortaleza (UNIFOR)} \\
Fortaleza, Cear\'a, Brasil \\
gabriel@email.com,
vlinsgurgel@gmail.com,
lorenzo@email.com,
davi@email.com}}

\maketitle

\begin{abstract}
A retinopatia diab\'etica \'e uma importante causa de perda de vis\~ao,
tornando relevante o desenvolvimento de m\'etodos automatizados para
apoiar sua triagem e diagn\'ostico. Este trabalho avaliou diferentes
arquiteturas de aprendizado profundo para a classifica\c{c}\~ao de
imagens de retina em cinco n\'iveis de severidade da doen\c{c}a,
utilizando o dataset \textit{Diabetic Retinopathy (resized)}. Foram
comparados modelos baseados em redes neurais convolucionais e Vision
Transformers, empregando transfer\^encia de aprendizado, fine-tuning,
aumento de dados e pondera\c{c}\~ao de classes para lidar com o
desbalanceamento da base. Os resultados indicaram que arquiteturas
modernas, como ConvNeXt-Small e EfficientNetV2-S, apresentaram os
melhores desempenhos gerais. O ConvNeXt-Small obteve o maior coeficiente
Kappa, enquanto o EfficientNetV2-S apresentou melhor equil\'ibrio entre
Kappa e F1 macro. Tamb\'em foi avaliada uma estrat\'egia de ensemble
baseada em \textit{Weighted Soft Voting}; entretanto, essa abordagem
n\~ao superou os modelos individuais na configura\c{c}\~ao testada. Os
resultados sugerem que a resolu\c{c}\~ao das imagens, o ajuste fino das
camadas e o tratamento do desbalanceamento s\~ao fatores determinantes
para o desempenho na classifica\c{c}\~ao multiclasse de retinopatia
diab\'etica.
\end{abstract}

\begin{IEEEkeywords}
Retinopatia diab\'etica, classifica\c{c}\~ao de imagens m\'edicas,
Deep Learning, CNNs, Vision Transformers, Transfer Learning.
\end{IEEEkeywords}

% =============================================================
\section{Introdu\c{c}\~ao}
% =============================================================

A retinopatia diab\'etica (RD) \'e uma complica\c{c}\~ao microvascular
do diabetes mellitus que afeta os vasos sangu\'ineos da retina e
constitui uma das principais causas de cegueira evit\'avel entre adultos
em idade produtiva. O crescimento cont\'inuo da popula\c{c}\~ao
diab\'etica mundial tem aumentado significativamente a demanda por
programas de rastreamento e detec\c{c}\~ao precoce da doen\c{c}a.
Segundo Grzybowski \cite{grzybowski2022}, a expans\~ao global do
diabetes e o aumento da incid\^encia de retinopatia diab\'etica tornam
cada vez mais necess\'aria a ado\c{c}\~ao de solu\c{c}\~oes escal\'aveis
para triagem oftalm\'ologica, especialmente em regi\~oes com recursos
limitados e escassez de especialistas.

Tradicionalmente, a identifica\c{c}\~ao da retinopatia diab\'etica \'e
realizada por oftalmologistas por meio da an\'alise de imagens de
retinografia. Entretanto, a grande quantidade de exames gerados em
programas de rastreamento populacional torna esse processo demorado,
oneroso e dependente da disponibilidade de profissionais qualificados.
Nesse contexto, t\'ecnicas de intelig\^encia artificial, especialmente
aquelas baseadas em aprendizado profundo, t\^em demonstrado grande
potencial para automatizar a an\'alise de imagens m\'edicas e auxiliar
a tomada de decis\~ao cl\'inica.

Entre os trabalhos pioneiros da \'area, Gulshan et al. \cite{gulshan2016}
desenvolveram um sistema baseado em redes neurais convolucionais profundas
para detec\c{c}\~ao autom\'atica de retinopatia diab\'etica em fotografias
de retina. O modelo alcan\c{c}ou elevados \'indices de sensibilidade e
especificidade em bases de valida\c{c}\~ao independentes, demonstrando
que algoritmos de aprendizado profundo podem atingir desempenho
compar\'avel ao de especialistas em tarefas de triagem da doen\c{c}a.

Nos \'ultimos anos, avan\c{c}os arquiteturais t\^em impulsionado ainda
mais o desempenho desses sistemas. Al\'em das tradicionais Redes Neurais
Convolucionais (Convolutional Neural Networks -- CNNs), arquiteturas
baseadas em mecanismos de aten\c{c}\~ao, conhecidas como Vision
Transformers (ViTs), passaram a apresentar resultados competitivos em
diversas tarefas de vis\~ao computacional. Em um estudo comparativo
envolvendo diferentes doen\c{c}as da retina, Sayg{\i}l{\i} e At{\i}c{\i}
\cite{saygili2025} observaram que o modelo Swin Transformer superou
arquiteturas CNN convencionais e outros modelos Transformer avaliados,
alcan\c{c}ando a maior acur\'acia entre os m\'etodos analisados.

Paralelamente, abordagens que combinam m\'ultiplos modelos tamb\'em v\^em
sendo investigadas com o objetivo de aumentar a robustez das
predi\c{c}\~oes. Rajeshwar et al. \cite{rajeshwar2025} propuseram uma
arquitetura h\'ibrida baseada em extra\c{c}\~ao de caracter\'isticas por
aprendizado profundo e ensemble por stacking, obtendo resultados
expressivos na detec\c{c}\~ao de retinopatia diab\'etica. Esses estudos
evidenciam o interesse crescente da comunidade cient\'ifica em explorar
diferentes arquiteturas e estrat\'egias de combina\c{c}\~ao de modelos
para melhorar o desempenho dos sistemas automatizados de apoio \`a
triagem oftalm\'ologica.

Apesar dos avan\c{c}os observados na literatura, a classifica\c{c}\~ao
autom\'atica da retinopatia diab\'etica ainda apresenta desafios
relevantes. Um deles \'e o desbalanceamento entre as classes, uma vez
que bases de dados p\'ublicas frequentemente possuem predomin\^ancia de
imagens sem sinais da doen\c{c}a em compara\c{c}\~ao com os n\'iveis
mais severos. Al\'em disso, trata-se de uma tarefa multiclasse e
ordinal, pois os n\'iveis de severidade seguem uma progress\~ao cl\'inica
de 0 a 4, como pode ser observado na Fig.~\ref{DR_levels}. Dessa forma,
m\'etricas globais como acur\'acia podem n\~ao refletir adequadamente o
desempenho do modelo em classes minorit\'arias, tornando necess\'aria a
an\'alise de m\'etricas como F1 macro e coeficiente Kappa ponderado.

\begin{figure*}[htbp]
\centering
\includegraphics[width=0.9\textwidth]{DR_levels.png}
\caption{Imagens do dataset por n\'ivel de severidade.}
\label{DR_levels}
\end{figure*}

Diante desse contexto, este trabalho tem como objetivo avaliar e comparar
diferentes arquiteturas de aprendizado profundo para a classifica\c{c}\~ao
de retinopatia diab\'etica em imagens do fundo da retina. Foram
investigados modelos baseados em CNNs e Vision Transformers, incluindo
ResNet50, DenseNet121, DenseNet169, EfficientNet-B0, EfficientNet-B2,
EfficientNetV2-S, ConvNeXt-Small e Swin-Tiny. Tamb\'em foram avaliadas
diferentes estrat\'egias de fine-tuning, aumento de dados, balanceamento
por pesos de classe e ensemble, buscando identificar as configura\c{c}\~oes
mais adequadas para a tarefa proposta.

% =============================================================
\section{Metodologia}
% =============================================================

A metodologia adotada consistiu na compara\c{c}\~ao experimental de
diferentes arquiteturas de aprendizado profundo para a classifica\c{c}\~ao
multiclasse de retinopatia diab\'etica em imagens de fundo de retina.
O fluxo experimental envolveu a prepara\c{c}\~ao da base de dados, o
pr\'e-processamento das imagens, o treinamento de modelos pr\'e-treinados
por transfer\^encia de aprendizado, a aplica\c{c}\~ao de estrat\'egias
de balanceamento e a avalia\c{c}\~ao dos modelos por m\'etricas
adequadas a um problema ordinal e desbalanceado.

\subsection{Dataset}

Foi utilizado o dataset p\'ublico \textit{Diabetic Retinopathy
(resized)}\cite{dataset}, derivado da competi\c{c}\~ao de retinopatia
diab\'etica do Kaggle \cite{APTOS2019}. A base cont\'em 35.126 imagens
de fundo de retina, rotuladas em cinco n\'iveis de severidade: classe 0,
sem retinopatia diab\'etica; classe 1, retinopatia leve; classe 2,
retinopatia moderada; classe 3, retinopatia severa; e classe 4,
retinopatia proliferativa. Como esses r\'otulos representam n\'iveis
progressivos de gravidade cl\'inica, o problema foi tratado como uma
tarefa de classifica\c{c}\~ao multiclasse ordinal.

Os metadados originais da base s\~ao disponibilizados em um arquivo CSV
contendo o identificador da imagem e o respectivo n\'ivel de severidade.
Para otimizar o treinamento dos modelos, as imagens foram reorganizadas
em diret\'orios separados por classe, utilizando uma estrutura compat\'ivel
com o carregamento por \textit{ImageFolder}. O conjunto foi dividido em
treino e valida\c{c}\~ao na propor\c{c}\~ao 80/20, com separa\c{c}\~ao
estratificada por classe e semente fixa igual a 42. N\~ao foi utilizado
um conjunto de teste externo independente.

\begin{table}[htbp]
\caption{Distribui\c{c}\~ao da base de dados por classe}
\label{tab:distribuicao_base}
\centering
\scriptsize
\begin{tabular}{c l r r r}
\hline
\textbf{Classe} & \textbf{Descri\c{c}\~ao} & \textbf{Total} &
\textbf{Treino} & \textbf{Val.} \\
\hline
0 & No DR         & 25.810 & 20.648 & 5.162 \\
1 & Mild          &  2.443 &  1.955 &   488 \\
2 & Moderate      &  5.292 &  4.234 & 1.058 \\
3 & Severe        &    873 &    699 &   174 \\
4 & Proliferative &    708 &    567 &   141 \\
\hline
\textbf{Total} & -- & \textbf{35.126} & \textbf{28.103} & \textbf{7.023} \\
\hline
\end{tabular}
\end{table}

Observa-se forte desbalanceamento entre as classes, com predomin\^ancia
da classe 0, correspondente a imagens sem sinais de retinopatia diab\'etica.
Esse fator motivou o uso de m\'etricas al\'em da acur\'acia e a aplica\c{c}\~ao
de estrat\'egias de pondera\c{c}\~ao durante o treinamento.

\subsection{Pr\'e-processamento e balanceamento dos dados}

As imagens foram redimensionadas de acordo com a configura\c{c}\~ao
experimental de cada modelo, sendo avaliadas resolu\c{c}\~oes de
224$\times$224, 256$\times$256, 384$\times$384, 512$\times$512 e
1024$\times$1024 pixels. Tamb\'em foi aplicada normaliza\c{c}\~ao com
m\'edia e desvio padr\~ao do ImageNet, uma vez que os modelos utilizados
foram inicializados com pesos pr\'e-treinados nesse conjunto.

Para aumentar a variabilidade dos dados de treinamento e reduzir o risco
de sobreajuste, foram aplicadas t\'ecnicas de \textit{data augmentation},
incluindo invers\~oes horizontal e vertical, rota\c{c}\~oes de at\'e
15 graus e altera\c{c}\~oes de cor por meio de \textit{ColorJitter}.
Essas transforma\c{c}\~oes foram aplicadas apenas ao conjunto de
treinamento, preservando o conjunto de valida\c{c}\~ao sem altera\c{c}\~oes
artificiais al\'em do redimensionamento e da normaliza\c{c}\~ao.

Devido ao desbalanceamento observado entre as classes, foi utilizada a
fun\c{c}\~ao de perda \textit{CrossEntropyLoss} ponderada por pesos de
classe, calculados a partir da distribui\c{c}\~ao dos r\'otulos no
conjunto de treinamento. Estrat\'egias adicionais, como
\textit{WeightedRandomSampler} e \textit{Focal Loss}, tamb\'em foram
exploradas, mas os principais experimentos utilizaram predominantemente
a pondera\c{c}\~ao da fun\c{c}\~ao de perda.

\subsection{Arquiteturas avaliadas e varia\c{c}\~oes experimentais}

Foram avaliadas arquiteturas baseadas em Redes Neurais Convolucionais e
Vision Transformers. Entre os modelos convolucionais, foram testadas as
arquiteturas EfficientNet-B0, EfficientNet-B2, EfficientNetV2-S,
ResNet50, DenseNet121, DenseNet169 e ConvNeXt-Small. Como representante
baseado em mecanismos de aten\c{c}\~ao visual, foi avaliado o modelo
Swin-Tiny.

Todos os modelos foram utilizados com pesos pr\'e-treinados em ImageNet,
caracterizando uma abordagem de transfer\^encia de aprendizado. A camada
final de classifica\c{c}\~ao foi substitu\'ida para produzir cinco
sa\'idas, correspondentes aos n\'iveis de severidade da retinopatia
diab\'etica. As varia\c{c}\~oes experimentais envolveram diferentes
resolu\c{c}\~oes de entrada, graus de congelamento do \textit{backbone},
otimizadores, taxas de aprendizado, tamanhos de lote e estrat\'egias de
fine-tuning.

Inicialmente, foram avaliadas arquiteturas com maior n\'ivel de
congelamento das camadas pr\'e-treinadas, funcionando como
configura\c{c}\~oes de refer\^encia. Em seguida, os modelos com
resultados mais promissores foram submetidos a ajustes mais espec\'ificos,
incluindo descongelamento parcial ou total do \textit{backbone}, aumento
da resolu\c{c}\~ao de entrada e altera\c{c}\~ao do otimizador. Essa
estrat\'egia permitiu analisar o impacto do fine-tuning e da
resolu\c{c}\~ao das imagens no desempenho dos classificadores.

\subsection{Estrat\'egia de treinamento}

Os experimentos foram implementados em Python utilizando as bibliotecas
PyTorch e torchvision, executados em GPUs NVIDIA compat\'iveis com CUDA.
Em raz\~ao das limita\c{c}\~oes de mem\'oria gr\'afica dispon\'ivel,
os tamanhos de lote foram ajustados conforme a arquitetura e a
resolu\c{c}\~ao das imagens utilizadas.

Foram utilizados os otimizadores Adam e AdamW, dependendo da
configura\c{c}\~ao experimental. Em parte dos experimentos, foram
aplicadas taxas de aprendizado distintas para o \textit{backbone} e
para a cabe\c{c}a classificadora, permitindo ajustes mais suaves nas
camadas pr\'e-treinadas e maior adapta\c{c}\~ao da camada final. Tamb\'em
foi utilizado o agendador \textit{ReduceLROnPlateau}, reduzindo a taxa de
aprendizado quando n\~ao havia melhora na m\'etrica de valida\c{c}\~ao.

O n\'umero m\'aximo de \'epocas variou entre 25 e 50, conforme o modelo
avaliado. Para evitar sobreajuste, foi empregado \textit{early stopping},
monitorando principalmente o coeficiente Kappa ponderado quadr\'atico no
conjunto de valida\c{c}\~ao. A vers\~ao final considerada para cada
modelo correspondeu ao melhor \textit{checkpoint} obtido durante o
treinamento, e n\~ao necessariamente \`a \'ultima \'epoca executada.

\subsection{Estrat\'egia de classifica\c{c}\~ao e ensemble}

A tarefa foi formulada como classifica\c{c}\~ao multiclasse direta. Para
cada imagem de entrada, o modelo produzia uma distribui\c{c}\~ao de
probabilidades sobre as cinco classes, sendo atribu\'ida como
predi\c{c}\~ao final a classe com maior probabilidade.

Al\'em da avalia\c{c}\~ao individual dos modelos, foi testada uma
estrat\'egia de ensemble baseada em \textit{Weighted Soft Voting}. Nessa
abordagem, as probabilidades previstas por diferentes modelos foram
combinadas por m\'edia ponderada, utilizando pesos proporcionais ao
desempenho individual de cada arquitetura no conjunto de valida\c{c}\~ao.
O ensemble foi avaliado com o objetivo de verificar se a combina\c{c}\~ao
de modelos de fam\'ilias distintas poderia melhorar a robustez das
predi\c{c}\~oes em compara\c{c}\~ao aos classificadores individuais.

\subsection{M\'etricas de avalia\c{c}\~ao}

O desempenho dos modelos foi avaliado por meio de acur\'acia,
precis\~ao macro, recall macro, F1-score macro e coeficiente Kappa
ponderado quadr\'atico. A acur\'acia foi utilizada como medida geral da
taxa de acerto; entretanto, devido ao forte desbalanceamento entre as
classes, essa m\'etrica n\~ao foi considerada suficiente de forma
isolada.

O F1-score macro foi utilizado por atribuir o mesmo peso a todas as
classes, independentemente da frequ\^encia de cada uma na base, o que
permite avaliar melhor o desempenho dos modelos nas classes
minorit\'arias. O coeficiente Kappa ponderado quadr\'atico foi adotado
como m\'etrica principal de compara\c{c}\~ao, pois considera a natureza
ordinal do problema, penalizando de forma mais intensa erros entre
classes distantes na escala de severidade. Tamb\'em foram analisadas
matrizes de confus\~ao para identificar os principais padr\~oes de erro
entre os n\'iveis de retinopatia diab\'etica.

% =============================================================
\section{Resultados e Discuss\~ao}
% =============================================================

\subsection{Compara\c{c}\~ao geral entre os modelos}

A Tabela~\ref{tab:comparacao_modelos} re\'une os principais resultados
obtidos. Para n\~ao sobrecarregar a compara\c{c}\~ao com todos os
experimentos realizados, foram inclu\'idos o baseline, o melhor
representante de cada fam\'ilia arquitetural e o ensemble final.

\begin{table}[htbp]
\caption{Compara\c{c}\~ao dos principais modelos avaliados}
\label{tab:comparacao_modelos}
\centering
\scriptsize
\begin{tabular}{l c c c c}
\hline
\textbf{Modelo} & \textbf{Res.} & \textbf{Accuracy} & \textbf{F1 Macro} & \textbf{Kappa} \\
\hline
EfficientNet-B0 (Baseline) & 224  & 0,74 & 0,27 & 0,26 \\
Swin-Tiny                  & 224  & 0,69 & 0,52 & 0,63 \\
DenseNet169                & 512  & 0,75 & 0,55 & 0,71 \\
EfficientNetV2-S           & 512  & 0,80 & 0,56 & 0,73 \\
ConvNeXt-Small             & 1024 & 0,82 & 0,53 & 0,75 \\
Ensemble                   & --   & 0,73 & 0,39 & 0,47 \\
\hline
\end{tabular}
\end{table}

O caso do EfficientNet-B0 \'e ilustrativo de um problema recorrente em
bases desbalanceadas: acur\'acia de 0,74 com kappa de apenas 0,26
indica que o modelo aprendeu, na pr\'atica, a prever quase sempre a
classe \textit{No\,DR} sem real capacidade de distinguir os n\'iveis da
doen\c{c}a. Essa discrep\^ancia refor\c{c}a a import\^ancia de avaliar
os modelos pelo kappa quadr\'atico e pelo F1 macro, e n\~ao apenas pela
acur\'acia global.

\subsection{Melhores modelos individuais}

Os tr\^es modelos com melhor desempenho --- ConvNeXt-Small,
EfficientNetV2-S e DenseNet169 --- apresentaram perfis bastante
distintos, e essa diferen\c{c}a \'e relevante do ponto de vista
cl\'inico.

O ConvNeXt-Small obteve o maior kappa do estudo (0,7453), mas com uma
fragilidade importante: recall de apenas 0,05 para a classe
\textit{Mild} (F1~=~0,09). Na pr\'atica, isso significa que o modelo
quase n\~ao detecta o est\'agio inicial da retinopatia, que \'e
justamente onde a interven\c{c}\~ao precoce teria maior impacto. O bom
resultado geral acaba mascarado pela domin\^ancia da classe
\textit{No\,DR}, que corresponde a 73\% das amostras e teve recall de
0,93.

O EfficientNetV2-S em 512\,px mostrou o melhor equil\'ibrio geral entre
as cinco classes, com kappa de 0,7313 e F1 macro de 0,56. Os recalls de
0,42 para \textit{Severe} e 0,67 para \textit{Proliferative} indicam
que o modelo conseguiu aprender caracter\'isticas das classes mais
graves, mesmo com muito menos exemplos dispon\'iveis. A classe
\textit{Mild} permaneceu a mais dif\'icil (recall 0,34; F1 0,27) ---
um resultado esperado, considerando que essa classe tem cerca de 10
vezes menos amostras que \textit{No\,DR} e que suas altera\c{c}\~oes
vasculares s\~ao visualmente sutis, podendo se confundir facilmente com
retinas saud\'aveis.

O DenseNet169 em 512\,px fechou o grupo dos melhores com kappa 0,7114 e
F1 macro de 0,55. Sua converg\^encia mais lenta --- pico na
\'epoca\,18 --- sugere um treinamento mais gradual e est\'avel, e seu
ponto forte ficou no recall de \textit{Proliferative} (0,67), o melhor
entre todos os modelos testados. A Tabela~\ref{tab:metricas_classe}
detalha as m\'etricas por classe dos tr\^es modelos.

\begin{table}[htbp]
\caption{M\'etricas por classe --- melhores modelos individuais}
\label{tab:metricas_classe}
\centering
\scriptsize
\setlength{\tabcolsep}{3.5pt}
\begin{tabular}{l l c c c}
\hline
\textbf{Modelo} & \textbf{Classe} & \textbf{Prec.} & \textbf{Rec.} & \textbf{F1} \\
\hline
\multirow{5}{*}{\shortstack[l]{ConvNeXt-S \\ (1024\,px)}}
  & No DR    & 0,90 & 0,93 & 0,92 \\
  & Mild     & 0,41 & 0,05 & 0,09 \\
  & Moderate & 0,55 & 0,48 & 0,51 \\
  & Severe   & 0,52 & 0,28 & 0,36 \\
  & Prolif.  & 0,71 & 0,55 & 0,62 \\
\hline
\multirow{5}{*}{\shortstack[l]{EffNetV2-S \\ (512\,px)}}
  & No DR    & 0,90 & 0,85 & 0,88 \\
  & Mild     & 0,22 & 0,34 & 0,27 \\
  & Moderate & 0,59 & 0,60 & 0,60 \\
  & Severe   & 0,44 & 0,42 & 0,43 \\
  & Prolif.  & 0,57 & 0,67 & 0,61 \\
\hline
\multirow{5}{*}{\shortstack[l]{DenseNet169 \\ (512\,px)}}
  & No DR    & 0,89 & 0,85 & 0,87 \\
  & Mild     & 0,17 & 0,30 & 0,22 \\
  & Moderate & 0,62 & 0,51 & 0,56 \\
  & Severe   & 0,39 & 0,53 & 0,45 \\
  & Prolif.  & 0,67 & 0,67 & 0,67 \\
\hline
\end{tabular}
\end{table}

\subsection{Curvas de aprendizagem}

% Descomente este bloco quando a figura estiver pronta:
% \begin{figure}[htbp]
% \centering
% \includegraphics[width=\columnwidth]{curvas_aprendizado.png}
% \caption{Curvas de \textit{loss} e kappa quadr\'atico de
% valida\c{c}\~ao por \'epoca para os tr\^es melhores modelos.}
% \label{fig:curvas}
% \end{figure}

As curvas de converg\^encia mostram comportamentos bem diferentes entre
os modelos. O EfficientNetV2-S convergiu de forma est\'avel ao longo de
10 \'epocas, sem sinais evidentes de sobreajuste at\'e que o
\textit{early stopping} interrompeu o treinamento. O DenseNet169 levou
mais tempo para estabilizar, mas manteve progress\~ao consistente at\'e
a \'epoca\,18. J\'a o ConvNeXt-Small atingiu seu melhor resultado na
\'epoca\,3 --- um comportamento not\'avel. A combina\c{c}\~ao de
\textit{backbone} amplamente descongelado com imagens de 1024\,px parece
ter acelerado demais o aprendizado inicial, limitando a capacidade do
modelo de refinar suas representa\c{c}\~oes nas \'epocas seguintes.

\subsection{Matrizes de confus\~ao}

\begin{figure}[htbp]
\centering
\includegraphics[width=\columnwidth]{matriz_confusao_effnet512.png}
\caption{Matriz de confus\~ao do EfficientNetV2-S em 512\,px
(\textit{checkpoint} da \'epoca\,10) no conjunto de valida\c{c}\~ao.}
\label{fig:confusao}
\end{figure}

A an\'alise da matriz de confus\~ao do EfficientNetV2-S em 512\,px
confirma que a classe \textit{Mild} \'e o ponto cr\'itico do modelo.
Grande parte dos seus erros ocorre por confus\~ao com \textit{No\,DR}
e \textit{Moderate}, o que faz sentido: les\~oes vasculares leves s\~ao
visualmente pr\'oximas tanto de retinas sem doen\c{c}a quanto de
est\'agios moderados da patologia. A classe \textit{Severe} tamb\'em
apresentou confus\~oes relevantes, principalmente com \textit{Moderate}
e \textit{Proliferative}. Em contrapartida, \textit{No\,DR} e
\textit{Moderate} tiveram as diagonais mais fortes. Vale notar que,
comparado aos experimentos em 384\,px, o modelo em 512\,px come\c{c}ou
a prever \textit{Severe} e \textit{Proliferative} de forma consistente,
o que n\~ao acontecia nas resolu\c{c}\~oes menores.

\subsection{Resultado do ensemble}

O ensemble por \textit{Weighted Soft Voting} apresentou desempenho
inferior ao esperado. Combinando seis modelos com pesos proporcionais ao
kappa individual de cada um, o resultado foi kappa de 0,4691 e F1 macro
de 0,39 --- pior que qualquer um dos tr\^es melhores modelos isolados,
como mostra a Tabela~\ref{tab:ensemble}.

\begin{table}[htbp]
\caption{Melhores modelos individuais versus ensemble}
\label{tab:ensemble}
\centering
\scriptsize
\begin{tabular}{l c c}
\hline
\textbf{Modelo} & \textbf{F1 Macro} & \textbf{Kappa} \\
\hline
ConvNeXt-Small (1024\,px)  & 0,53 & 0,75 \\
EfficientNetV2-S (512\,px) & 0,56 & 0,73 \\
DenseNet169 (512\,px)      & 0,55 & 0,71 \\
\textbf{Ensemble}          & \textbf{0,39} & \textbf{0,47} \\
\hline
\end{tabular}
\end{table}

A hip\'otese mais prov\'avel para esse resultado \'e a incompatibilidade
entre as distribui\c{c}\~oes de probabilidade dos modelos: quando as
sa\'idas n\~ao est\~ao calibradas numa mesma escala, a m\'edia ponderada
tende a produzir predi\c{c}\~oes inconsistentes. Al\'em disso, os pesos
do ensemble foram definidos de forma fixa, sem otimiza\c{c}\~ao sobre o
conjunto de valida\c{c}\~ao. N\~ao se pode descartar tamb\'em algum
problema no carregamento dos \textit{checkpoints} de arquiteturas
distintas durante a execu\c{c}\~ao. Assim, na configura\c{c}\~ao
avaliada, o ensemble n\~ao apresentou ganho suficiente para justificar
seu custo computacional adicional.

\subsection{Discuss\~ao dos principais achados}

Um dos achados mais relevantes ao longo dos experimentos foi o impacto
da resolu\c{c}\~ao. Aumentar o EfficientNetV2-S de 384\,px para 512\,px
elevou o kappa de 0,57 para 0,73 e quase dobrou o recall de
\textit{Mild}, de 0,26 para 0,34. A explica\c{c}\~ao \'e direta:
microaneurismas e exsudatos caracter\'isticos dos est\'agios iniciais da
retinopatia s\~ao estruturas pequenas, e resolu\c{c}\~oes menores podem
reduzir ou eliminar detalhes importantes que o modelo precisaria aprender
a reconhecer.

O desbalanceamento da base, por outro lado, provou ser mais resistente
do que o esperado. A pondera\c{c}\~ao de classes na fun\c{c}\~ao de
perda ajudou, mas n\~ao foi suficiente para compensar uma raz\~ao de
36:1 entre \textit{No\,DR} e \textit{Proliferative}. A situa\c{c}\~ao
da classe \textit{Mild} \'e ainda mais delicada porque combina dois
problemas ao mesmo tempo: poucos exemplos e alta similaridade visual
com outras classes. O ConvNeXt-Small ilustra bem o risco de n\~ao
observar as m\'etricas por classe: kappa de 0,7453, mas F1 de 0,09
para \textit{Mild}. Um modelo assim poderia passar por uma
avalia\c{c}\~ao superficial e ainda assim ser inapropriado para uso
cl\'inico real.

Quanto ao ensemble, os resultados indicam que combinar modelos sem
calibra\c{c}\~ao pr\'evia pode ser contraproducente. Sem garantir que
as distribui\c{c}\~oes de probabilidade de cada modelo estejam numa
mesma escala, a combina\c{c}\~ao por m\'edia ponderada n\~ao agrega
informa\c{c}\~ao de forma coerente. \'E uma li\c{c}\~ao importante:
ensemble n\~ao \'e uma solu\c{c}\~ao autom\'atica para melhorar
desempenho --- a implementa\c{c}\~ao importa tanto quanto a estrat\'egia.

% =============================================================
\section{Conclus\~ao}
% =============================================================

Este trabalho avaliou diferentes configura\c{c}\~oes de arquiteturas de
aprendizado profundo --- incluindo CNNs modernas e um Vision Transformer
--- para a classifica\c{c}\~ao de cinco n\'iveis de severidade de
retinopatia diab\'etica em 35.126 imagens de fundo de retina. Foram
investigados o impacto da resolu\c{c}\~ao de entrada, do fine-tuning
parcial do \textit{backbone}, da pondera\c{c}\~ao de classes e da
combina\c{c}\~ao de modelos por ensemble.

O EfficientNetV2-S em 512$\times$512 pixels apresentou o melhor
equil\'ibrio cl\'inico, com kappa de 0,7313 e F1 macro de 0,56, sendo
o modelo com maior capacidade de detectar os cinco n\'iveis de forma
consistente. O ConvNeXt-Small alcan\c{c}ou o maior kappa absoluto
(0,7453), mas seu recall de 0,05 para \textit{Mild} compromete sua
utilidade em cen\'arios de triagem precoce. O ensemble n\~ao superou os
modelos individuais, evidenciando que a combina\c{c}\~ao de modelos
exige calibra\c{c}\~ao adequada para ser efetiva.

De forma geral, os experimentos mostraram que resolu\c{c}\~ao e
fine-tuning s\~ao fatores decisivos para o desempenho, mas que o
desbalanceamento extremo da base permanece o principal obst\'aculo.
Nenhuma das estrat\'egias testadas foi suficiente para garantir boa
detec\c{c}\~ao em todas as classes simultaneamente --- especialmente em
\textit{Mild}, onde a dificuldade \'e tanto estat\'istica quanto visual.

O estudo tem limita\c{c}\~oes relevantes: n\~ao foi utilizado conjunto
de teste externo independente, as restri\c{c}\~oes de mem\'oria de GPU
impediram explorar 1024\,px em todas as arquiteturas, e os experimentos
com \textit{focal loss} e \textit{sampler} n\~ao foram conclu\'idos.
Para trabalhos futuros, as dire\c{c}\~oes mais promissoras incluem
\textit{data augmentation} seletiva nas classes minorit\'arias ---
aplicando transforma\c{c}\~oes mais intensas especificamente sobre
\textit{Mild}, \textit{Severe} e \textit{Proliferative} como forma de
oversampling sem introduzir ru\'ido nas classes j\'a bem representadas
---, valida\c{c}\~ao cruzada estratificada, \textit{threshold moving}
por classe, calibra\c{c}\~ao de probabilidades para viabilizar o ensemble
e aplica\c{c}\~ao de Grad-CAM para interpretar as regi\~oes de interesse
identificadas pelos modelos.

% =============================================================
\begin{thebibliography}{00}
% =============================================================
\bibitem{grzybowski2022} A. Grzybowski, ``Artificial intelligence for
diabetic retinopathy screening,'' \textit{Acta Ophthalmol.}, vol.~100,
2022, doi: 10.1111/j.1755-3768.2022.15371.

\bibitem{gulshan2016} V. Gulshan et al., ``Development and Validation of
a Deep Learning Algorithm for Detection of Diabetic Retinopathy in
Retinal Fundus Photographs,'' \textit{JAMA}, vol.~316, no.~22,
pp.~2402--2410, Dec. 2016, doi: 10.1001/jama.2016.17216.

\bibitem{saygili2025} A. Sayg{\i}l{\i} and \"O.\,M. At{\i}c{\i},
``Enhancing Retinal Disease Detection With the Swin Transformer: A
Comprehensive Comparative Analysis,''
\textit{Concurrency Computat. Pract. Exper.}, vol.~37, no.~27--28,
p.~e70366, 2025, doi: 10.1002/cpe.70366.

\bibitem{rajeshwar2025} S. Rajeshwar, S. Thaplyal, A.\,M., and
G. Siva Shanmugam, ``Diabetic Retinopathy Detection Using DL-Based
Feature Extraction and a Hybrid Attention-Based Stacking Ensemble,''
\textit{Adv. Public Health}, vol.~2025, p.~8863096, 2025,
doi: 10.1155/adph/8863096.

\bibitem{dataset} ilovescience, ``Diabetic Retinopathy Resized,''
Kaggle, 2019. Dispon\'ivel em:
https://www.kaggle.com/datasets/tanlikesmath/diabetic-retinopathy-resized.
Acesso em: 4 jun. 2026.

\bibitem{APTOS2019} Karthik, Maggie, e S. Dane, ``APTOS 2019 Blindness
Detection,'' Kaggle, 2019. Dispon\'ivel em:
https://kaggle.com/competitions/aptos2019-blindness-detection.
Acesso em: 4 jun. 2026.
\end{thebibliography}

\end{document}